\documentclass[12pt]{article}
\usepackage{palestra1}
\begin{document}
\begin{ficha}{Fuerza y bando}{El número y su signo}{1}

\begin{encargo}
Cada tirador lleva un número. El \textbf{signo} dice su bando:
\textcolor{izq}{\textbf{−}}, rojo, tira hacia la izquierda;
\textcolor{der}{\textbf{+}}, azul, hacia la derecha. El número sin el signo es su
\textbf{fuerza}.
\end{encargo}

\bloque{1}{Mira cómo se hace}{En el campo}

\vspace{1.5mm}
\begin{minipage}[c]{.68\linewidth}
\paso{1}\textbf{¿De qué color es?}\ Rojo: tira hacia la izquierda, \textcolor{izq}{\textbf{−}}.\\[1.5mm]
\paso{2}\textbf{¿Qué número lleva en la cabeza?}\ Un 4: su fuerza es 4.\\[1.5mm]
\paso{3}\textbf{Su número, con su signo:}\ \nm{-4}.
\end{minipage}%
\begin{minipage}[c]{.32\linewidth}\centering
\tirador[1.1]{-4}
\end{minipage}

\vspace{1.5mm}
\begin{listillon}
La fuerza es el \textbf{valor absoluto}: el número sin su signo. Se escribe
entre barras: $|{-4}| = 4$.
\end{listillon}

\bloque{2}{Ahora tú}{En el campo}\quad{\small su número con signo y su fuerza}

\vspace{1.5mm}
{\renewcommand{\arraystretch}{1.6}
\begin{tabular}{@{}r@{\hspace{4mm}}*{4}{c@{\hspace{10mm}}}@{}}
& \tirador[.85]{3} & \tirador[.85]{-2} & \tirador[.85]{-6} & \tirador[.85]{5}\\
& \ej{a} & \ej{b} & \ej{c} & \ej{d}\\
número & \casillita[8mm] & \casillita[8mm] & \casillita[8mm] & \casillita[8mm]\\
fuerza & \casillita[8mm] & \casillita[8mm] & \casillita[8mm] & \casillita[8mm]\\
\end{tabular}}

\bloque{3}{En la regla}{Con la regla}\quad{\small más a la izquierda es menor}

\vspace{2mm}
\reglavacia[6mm]{8}

\vspace{2mm}
Márcalos en la regla y pon el signo $<$ o $>$:\\[2mm]
\ej{a}$-3$ \signo\ $+2$\qquad\qquad \ej{b}$-5$ \signo\ $-1$\qquad\qquad \ej{c}$+4$ \signo\ $-6$

\alto{$-5$ tira más fuerte que $-1$, pero $-5$ es \textbf{menor}: está más a
la izquierda en la regla.}

\reto{¿Puede haber un tirador de fuerza 0? ¿En qué bando tiraría?}

\comprueba{Cada código abre ese duelo en Practicar.}{%
\qr{3a}{j=fuerza\&a=-3\&b=2\&q=menor}\hfill\qr{3b}{j=fuerza\&a=-5\&b=-1\&q=menor}\hfill
\qr{3c}{j=fuerza\&a=4\&b=-6\&q=mayor}}

\end{ficha}

% ─────────────────────────────── REVERSO ──────────────────────────────────
\begin{dorso}{Fuerza y bando}{Más duelos}{1}

\bloque{1}{¿Quién tira más fuerte?}{Sobre el papel}\quad{\small rodéalo}

\vspace{2mm}
{\renewcommand{\arraystretch}{2}
\begin{tabular}{@{}l@{\hspace{16mm}}l@{}}
\ej{a}\nm{-6}\quad o\quad \nm{4} & \ej{b}\nm{2}\quad o\quad \nm{-7}\\
\ej{c}\nm{-3}\quad o\quad \nm{-8} & \ej{d}\nm{5}\quad o\quad \nm{-1}\\
\end{tabular}}

\bloque{2}{¿Quién es mayor?}{Con la regla}\quad{\small márcalos y pon el signo}

\vspace{2mm}
\reglavacia[6mm]{8}

\vspace{2mm}
{\renewcommand{\arraystretch}{2}
\begin{tabular}{@{}l@{\hspace{16mm}}l@{}}
\ej{a}$-4$ \signo\ $+1$ & \ej{b}$-2$ \signo\ $-6$\\
\ej{c}$0$ \signo\ $-3$ & \ej{d}$+3$ \signo\ $-5$\\
\end{tabular}}

\bloque{3}{Problemas}{Sobre el papel}

\vspace{1.5mm}
\ej{a}En Burgos hay $-3$\,°C y en Madrid, $+2$\,°C. ¿Dónde hace más frío?\hfill\hueco[2.4cm]

\vspace{4mm}
\ej{b}El garaje está en la planta $-2$ y el trastero en la $-4$. ¿Cuál está más
abajo?\hfill\hueco[2.4cm]

\vspace{4mm}
\begin{semaforo}
\sfila{Escribo el número de un tirador con su signo}
\sfila{Digo la fuerza de un número}
\sfila{Sé cuál es mayor mirando la regla}
\end{semaforo}

\comprueba{Cada código abre ese duelo en Practicar.}{%
\qr{1a}{j=fuerza\&a=-6\&b=4\&q=fuerza}\hfill\qr{1c}{j=fuerza\&a=-3\&b=-8\&q=fuerza}\hfill
\qr{2a}{j=fuerza\&a=-4\&b=1\&q=mayor}\hfill\qr{2b}{j=fuerza\&a=-2\&b=-6\&q=mayor}\hfill
\qr{2d}{j=fuerza\&a=3\&b=-5\&q=mayor}}
\end{dorso}
\end{document}
